\documentclass[conference]{IEEEtran}
\usepackage{cite}
\usepackage{amsmath,amssymb,amsfonts}
\usepackage{graphicx}
\usepackage{booktabs}
\usepackage{url}
\usepackage{hyperref}

\begin{document}

\title{Sensor Fault Discovery in Automatic Weather Station Networks:
Separating Instrument Failures from Genuine Extremes Without Labels}

\author{\IEEEauthorblockN{Jeetaksh Luhaniwal}
\IEEEauthorblockA{Department of Computer Science / AI \& ML\\
INT-396 Unsupervised Learning}}

\maketitle

\begin{abstract}
Automatic weather station networks feed forecasting and disaster-response systems, where a silently
faulty sensor corrupts every downstream decision. We present an unsupervised pipeline that separates
\emph{failing instruments} from \emph{genuine environmental extremes} using only unlabelled
multi-station data. Eight Indian stations were sampled hourly from ERA5 reanalysis for 2023 and
augmented with six injected sensor faults and four physically coherent extreme-event types. Five
models implemented from scratch --- k-Means++, DBSCAN, PCA reconstruction error, Isolation Forest and
Local Outlier Factor --- vote per reading, and a consistency rule tree separates faults from
extremes. Against a fixed-range quality-control baseline, the detector raises fault detection rate
from 0.281 to 0.393 and F1 from 0.439 to 0.451, with precision@50 of 0.96, while keeping the
extreme false-flag rate at 0.85\%. Ground truth was used only for scoring, never for fitting.
\end{abstract}

\begin{IEEEkeywords}
unsupervised learning, anomaly detection, weather stations, sensor faults, clustering,
isolation forest, PCA.
\end{IEEEkeywords}

\section{Introduction}
An automatic weather station that fails silently is more dangerous than one that fails loudly. A
range check catches a thermometer reading 200~$^\circ$C. It does not catch a sensor stuck at
24.0~$^\circ$C for eighteen hours, nor one whose calibration drifts by three degrees over a week,
nor one whose resolution degrades to whole degrees. All three produce plausible numbers.

The same property makes the problem hard in the other direction: a genuine heat wave also produces
numbers outside the comfortable range. An operator who silences every out-of-range alert will silence
the events that matter most.

This work asks two questions with one detector. Which readings are anomalous, given no labels? And
for each anomaly, is the instrument failing or the weather extreme?

Our contributions are: (i) a consensus of five classical unsupervised detectors operating on
physically meaningful features; (ii) a transparent discriminator that uses cross-sensor coherence,
neighbour agreement and physical plausibility; (iii) an evaluation that reports performance on faults
and on extremes separately, against a mandatory baseline.

\section{Background}
Unsupervised methods separate structure from noise without labels. Partition-based clustering
minimises within-cluster variance; density-based clustering marks low-density points as noise by
construction; reconstruction methods learn the data manifold and measure distance from it; tree-based
isolation measures how few random splits are needed to isolate a point.

Three distinctions matter here. \emph{Anomaly} detection finds rare points relative to the observed
distribution. \emph{Novelty} detection finds points unlike a training distribution. We do the former
during scoring and use a clean reference period for threshold calibration, which leans on the latter.
Distance choice also changes the meaning of ``similar'': Euclidean is magnitude-sensitive, Manhattan
resists single-coordinate spikes, cosine compares shape rather than scale.

\section{Data and Problem Formulation}
Hourly readings for eight stations across 2023 were taken from the Open-Meteo historical API, which
serves ERA5 reanalysis under CC~BY~4.0~\cite{openmeteo,era5}. The tidy table holds 280{,}320 rows of
(station, hour, sensor, value). Sensors are temperature, humidity, wind speed and pressure;
precipitation is retained as an auxiliary signal.

ERA5 is quality-controlled and contains no failing instruments. Labels were therefore produced by
injection: six fault types applied to a single sensor of a single station per episode, and four
extreme types applied coherently across all sensors of a station. Faults cover 1.80\% of readings and
extremes 4.30\%. Label columns exist only in the evaluation harness.

\begin{table}[ht]
\centering
\caption{Injected fault and extreme taxonomy.}
\label{tab:taxonomy}
\begin{tabular}{lll}
\toprule
Type & Mechanism & Duration \\
\midrule
bias & offset 3--8$\sigma$ & 12--72\,h \\
drift & ramp 2--5$\sigma$ & 3--10\,d \\
stuck & frozen value & 6--48\,h \\
spike & jumps 5--12$\sigma$ & samples \\
precision & quantisation & 24--120\,h \\
dropout & missing values & 3--24\,h \\
\midrule
heat wave & $+6$--$12\,^\circ$C, drier & 3--7\,d \\
cold snap & $-6$--$15\,^\circ$C & 2--5\,d \\
storm & wind $\times2$--4, pressure drop & 12--36\,h \\
heavy rain & humidity $\to$ 90--100\% & 6--24\,h \\
\bottomrule
\end{tabular}
\end{table}

\section{Methodology}
\subsection{Features}
Ten features are computed per reading from measured values only: standardised value, deviation from a
24-hour median, rate of change, 24-hour rolling mean and standard deviation, a flatline indicator for
six identical consecutive readings, cyclic hour encoding, an association-rule violation count, a
physical-rule violation count, and a missing-value indicator. Scaling statistics come from the first
30\% of the timeline, computed per station so that Delhi and Shimla are not forced onto a shared
scale.

\subsection{Consensus scoring}
Each of the five models produces one score. Scores are converted to votes against their own 99th
percentile on the reference period, and two or more votes mark a reading anomalous. The soft score is
the vote fraction.

\subsection{Discrimination}
Consecutive anomalous hours for one station and sensor form an episode. Three signals are computed
per episode: the number of sensors flagged at the same station-hour, the share of the other seven
stations whose same sensor deviates, and whether any hard physical rule is violated. The decisive
signal proved to be the coherence ratio --- deviation of the \emph{other} sensors divided by deviation
of the flagged one. Injected faults move one sensor (ratio 0.11 median at episode level) while
coherent extremes move several (ratio 0.34).

\subsection{Baseline}
Fixed limits per sensor (temperature $-5$ to 45~$^\circ$C, humidity 0--100\%, wind 0--120\,km/h,
pressure 950--1050\,hPa) flag anything outside, with no context and no learning.

\section{Experimental Setup}
Seed 42, Python 3.10, NumPy and pandas only in the model path; scikit-learn is absent from production
code and used at most as a numerical oracle during development. The detector runs in about 100\,s and
the discriminator in about 2\,s on 280{,}320 readings.

\section{Results}
\begin{table}[ht]
\centering
\caption{Detector versus fixed-range baseline.}
\label{tab:main}
\begin{tabular}{lccc}
\toprule
Metric & Consensus & Range check \\
\midrule
Fault detection rate $\uparrow$ & 0.393 & 0.281 \\
Fault classification rate $\uparrow$ & 0.274 & 0.281 \\
Extreme false-flag rate $\downarrow$ & 0.0085 & 0.0037 \\
Precision (fault) $\uparrow$ & 0.530 & 1.000 \\
F1 (fault) $\uparrow$ & 0.451 & 0.439 \\
False positives on clean $\downarrow$ & 0.0056 & 0.0000 \\
\bottomrule
\end{tabular}
\end{table}

\begin{table}[ht]
\centering
\caption{Precision at top-$K$ by anomaly score.}
\label{tab:precision}
\begin{tabular}{lccc}
\toprule
Method & P@10 & P@25 & P@50 \\
\midrule
Consensus scoring & 1.00 & 1.00 & 0.96 \\
Range check & 1.00 & 1.00 & 1.00 \\
\bottomrule
\end{tabular}
\end{table}

Per-fault recall explains the gap in Table~\ref{tab:main}. Dropout reaches 1.00 and stuck 0.81,
because both are invisible to a range check by construction. Bias reaches 0.51 and spikes 0.63. Drift
(0.11) and precision degradation (0.04) remain weak: a slow ramp stays inside normal ranges, and
quantisation noise is subtle at hourly resolution.

The baseline's perfect precision is a consequence of its conservatism, not superiority: it flags only
0.28 of faults, so almost everything it flags is correct. It cannot rank, which is why the consensus
score matters operationally.

\section{Discussion and Limitations}
Faults and extremes are synthetic. Magnitudes were authored, so these numbers are optimistic relative
to a real deployment, and ERA5 is already quality-controlled. The rule tree encodes meteorological
reasoning and would misclassify a novel fault that happens to move several sensors coherently.
Episodes fragment: a long event often produces several short episodes rather than one, because the
vote threshold is strict. Precision of 0.53 means roughly two flags per true fault, which is an
acceptable trade for a monitoring alarm but must be stated.

\section{Conclusion}
A consensus of five classical unsupervised detectors, scored on physically meaningful features and
followed by a coherence-based discriminator, detects more sensor faults than a fixed range check
while keeping genuine extremes from being mislabelled as failures. The remaining weaknesses --- drift,
precision loss, and episode fragmentation --- are specific and addressable with sequence models and
longer horizon features.

\begin{thebibliography}{9}
\bibitem{openmeteo}
P. Zippenfenig, ``Open-Meteo.com Weather API,'' Zenodo, 2023.
doi: 10.5281/zenodo.7970649.

\bibitem{era5}
H. Hersbach et al., ``ERA5 hourly data on single levels from 1940 to present,'' ECMWF, 2023.
doi: 10.24381/cds.adbb2d47.

\bibitem{breifor}
F. T. Liu, K. M. Ting, and Z.-H. Zhou, ``Isolation Forest,'' in \emph{Proc. ICDM}, 2008, pp. 413--422.

\bibitem{lof}
M. M. Breunig et al., ``LOF: Identifying Density-Based Local Outliers,'' in \emph{Proc. ACM SIGMOD},
2000, pp. 93--104.

\bibitem{kmeanspp}
D. Arthur and S. Vassilvitskii, ``k-means++: The Advantages of Careful Seeding,'' in \emph{Proc.
SODA}, 2007, pp. 1027--1035.

\bibitem{dbscan}
M. Ester et al., ``A Density-Based Algorithm for Discovering Clusters in Large Spatial Databases with
Noise,'' in \emph{Proc. KDD}, 1996, pp. 226--231.

\bibitem{pca}
I. T. Jolliffe, \emph{Principal Component Analysis}, 2nd ed. Springer, 2002.

\bibitem{silhouette}
P. J. Rousseeuw, ``Silhouettes: A Graphical Aid to the Interpretation and Validation of Cluster
Analysis,'' \emph{J. Comput. Appl. Math.}, vol. 20, pp. 53--65, 1987.

\bibitem{apriori}
R. Agrawal and R. Srikant, ``Fast Algorithms for Mining Association Rules,'' in \emph{Proc. VLDB},
1994, pp. 487--499.

\bibitem{fpgrowth}
J. Han, J. Pei, and Y. Yin, ``Mining Frequent Patterns without Candidate Generation,'' in
\emph{Proc. SIGMOD}, 2000, pp. 1--12.
\end{thebibliography}

\end{document}
